\section{Experiments}\label{sec:experiments}
\subsection{Datasets and Few-Shot Splits}
We evaluate on MVTec AD~\cite{mvtec2019} and VisA~\cite{visa2022}. MVTec AD contains industrial object and texture categories with pixel-level defect masks. VisA contains a broader set of industrial categories and is used as the main full-scale corruption and calibration-shift benchmark. For each class, support sets are sampled from normal training images with fixed seeds. Unless otherwise specified, $k\in\{1,2,4,8\}$ is used for clean accuracy-storage studies, and $k\in\{4,8\}$ is used for the full conformal corruption study because LOIO conformal requires at least two support images and is most relevant in moderate few-shot settings.

\subsection{Evaluation Protocols}
We organize experiments into five protocol groups. First, the clean few-shot protocol evaluates image AUROC, AP, calibration, and storage on MVTec and VisA. Second, the accuracy-storage protocol compares PCA64 and PCA128 to test whether a small subspace increase improves accuracy without approaching memory-bank storage. Third, the corruption reliability protocol evaluates Gaussian noise, blur, brightness/contrast shift, and JPEG compression at $k\in\{4,8\}$ with at most 120 label-stratified images per case: on VisA across all 12 classes and five seeds (480 cases, 56{,}000 image rows), and on MVTec across all 15 classes and three seeds (360 cases, 37{,}464 image rows); VisA serves as the primary full-scale benchmark with the larger seed budget, while the MVTec replication uses three seeds for compute reasons, which its paired per-cell statistics account for. Fourth, representative routing protocols evaluate MVTec$\rightarrow$VisA, VisA$\rightarrow$MVTec, leave-one-class-out (LOCO), and within-dataset class splits. Fifth, the operational false-alarm protocol thresholds conformal p-values at nominal levels $\alpha\in\{0.01,0.05,0.10,0.20\}$ and reports false-alarm rate on normal images, detection power, alarm precision, and the attainable-alpha structure; the SC3R protocol additionally selects thresholds on held-out source-class normals with label-stratified random image sampling and a fixed sampling manifest, and is judged against a decision gate registered before the stratified run (nonzero power at $\alpha\le0.10$, mean false-alarm rate at most $\alpha{+}0.02$, no-harm at least 80\%, no target anomaly labels, hierarchical class--seed bootstrap support).

\subsection{Naming Convention}
CRR denotes the overall framework (decoupled ranking plus reliability layer). The detector instantiations are named by their subspace configuration: CRR/PCA64 (default, also called CalibSubspaceHead in configuration files) and CRR/PCA128. Reliability views are named by their route: scalar/Vector/Shift-Aware/weighted Platt, temperature scaling, isotonic regression, histogram binning, LOIO conformal, and weighted conformal. Tables use these names consistently.

\subsection{Baselines and Ablations}
We compare against a controlled DINOv2 nearest-neighbor memory-bank baseline, a SubspaceAD-style PCA residual baseline, scalar Platt, Vector Platt, Shift-Aware Platt, weighted Platt, temperature scaling, isotonic regression, histogram binning, LOIO conformal, and weighted conformal. The controlled memory-bank implementation is not an official reproduction of PatchCore or AnomalyDINO; it holds the frozen DINOv2 feature pipeline and data split fixed to isolate the scoring and storage trade-off. In addition, we run the official AnomalyDINO code (ViT-S/14 at 448px, agnostic preprocessing) on both MVTec and VisA (official 1-cls split) for $k\in\{1,4,8\}$ with three seeds and report it as a separate accuracy reference. Official SubspaceAD representative runs are included as a novelty guardrail. Their strong performance confirms that frozen DINOv2 subspace residuals are an existing competitive ranking mechanism; CRR should therefore be judged by reliability, storage, and diagnostic value rather than by claiming the residual itself as novel.

\subsection{Metrics}
We separate ranking metrics from reliability metrics. AUROC and AP are computed from the raw PCA/subspace residual score $s(x)$. Pixel-level localization is measured by pixel AUROC and AU-PRO when masks are available. Reliability is measured operationally by empirical false-alarm rate at fixed $\alpha$, detection power, alarm precision, discrete uniformity of normal p-values, and risk--coverage curves; expected calibration error (ECE), Brier score, and negative log-likelihood (NLL) are reported as secondary metrics using the selected reliability probability or probability-like conformal view. Storage is measured as retained reference/model state per class. Robustness experiments include corruption shift and a label-aware PCA-residual FGSM surrogate; the latter is used only as a fragility diagnostic.

\subsection{Reproducibility and Claim Rules}
All runs are indexed by dataset, class, $k$, seed, model variant, calibration mode, and corruption/attack setting. The core rule is that reliability routes cannot claim ranking improvement unless the raw score changes. Conversely, ranking methods cannot claim calibrated probabilities unless ECE/Brier/NLL improve. This separation prevents a common ambiguity in AD papers: conflating a better sorted list with a more trustworthy alarm probability.
